% 第4章 4.5.3：边界迹示意。(a) 有界Lipschitz区域 Omega，边界点附近的球内边界可写成 Lipschitz 函数图像 x_d=gamma(x')，
% 区域位于图像上方；(b) 平直化后的半空间 {x_d>0}：沿竖直线段从边界点 (x',0) 向内积分，|u(x',0)|^2<=∫(|u|^2+|∂_d u|^2)dx_d。
\begin{tikzpicture}[font=\figurefont\footnotesize,line cap=round,line join=round]
% (a)
\begin{scope}
\fill[FigureBlueFill] plot[smooth cycle,tension=.75] coordinates {(0,0) (1.4,-0.35) (2.9,0.05) (3.3,1.2) (2.6,2.2) (1.2,2.25) (0.1,1.6) (-0.3,0.7)};
\draw[FigureBlue,line width=1.0pt] plot[smooth cycle,tension=.75] coordinates {(0,0) (1.4,-0.35) (2.9,0.05) (3.3,1.2) (2.6,2.2) (1.2,2.25) (0.1,1.6) (-0.3,0.7)};
\node[text=FigureInk] at (1.6,1.2) {$\Omega$};
\node[text=FigureBlue,anchor=north] at (2.55,-0.25) {$\partial\Omega$};
\draw[FigureYellow,line width=.8pt,dashed] (1.4,-0.33) circle (0.62);
\fill[FigureInk] (1.4,-0.33) circle (1.3pt);
\node[text=FigureInk,anchor=north,font=\figurefont\scriptsize] at (1.4,-0.98) {$x_0$};
\node[font=\figurefont\footnotesize] at (1.5,2.75) {(a) 局部写成图像};
\end{scope}
\draw[-{Latex[length=2mm]},FigureStroke,line width=.7pt] (3.7,0.6) -- node[above,font=\FigureNoteFont\scriptsize,text=FigureMuted]{平直化} (4.7,0.6);
% (b)
\begin{scope}[xshift=5.1cm]
\fill[FigureBlueFill] (0,0) rectangle (3.4,2.1);
\draw[FigureBlue,line width=1.0pt] (0,0) -- (3.4,0);
\node[text=FigureBlue,anchor=north,font=\figurefont\scriptsize] at (2.9,-0.05) {$x_d=0$};
\node[text=FigureInk,font=\figurefont\scriptsize] at (2.75,1.7) {$x_d>0$};
\draw[FigureGreen,line width=1.1pt,-{Latex[length=1.8mm]}] (1.3,1.7) -- (1.3,0.05);
\fill[FigureInk] (1.3,0) circle (1.3pt);
\node[text=FigureInk,anchor=north,font=\figurefont\scriptsize] at (1.3,-0.05) {$(x',0)$};
\node[text=FigureGreen,anchor=west,font=\figurefont\scriptsize,align=left] at (1.4,0.95) {沿$x_d$积分};
\node[font=\figurefont\footnotesize] at (1.7,2.75) {(b) 半空间};
\end{scope}
\end{tikzpicture}
